% R10 manuscript fragment. Requires amsmath, amssymb, graphicx.
% See PROOF.md for full arguments, source links and trust boundaries.
% This fragment has not been compiled as part of the original manuscript.
\section{A finite swallowtail window in the original controls}

Fix the positive modified kernel from the preceding construction at
its exact order-four amplitude $\epsilon_4$. Let $G$ be its integral,
let $Q=(t_Q,\lambda_Q,\mu_Q)$, $\nu=0$, be the exact identified point,
and put
\[
 s=t-t_Q,\quad \ell=\lambda-\lambda_Q,\quad m=\mu-\mu_Q,
 \qquad g=24G/G_4(Q).
\]
The normalizing multiplier is a nonzero constant. The zeros and their
multiplicities are unchanged; at the exact origin
$g_0=g_1=g_2=g_3=0$ and $g_4=24$. Subscripts denote $s$ derivatives.
The heat identities are
\[
 \partial_\ell g_j=-g_{j+2}/4,\qquad
 \partial_m g_j=g_{j+4}/16,\qquad
 \partial_\nu g_j=-g_{j+6}/64.
\]

\paragraph{Theorem 1 (finite root count).}
Let $T=2^{-7}$ and
\[
 P=\{(\ell,m,\nu): |\ell|\leq2^{-17},\ |m|,|\nu|\leq2^{-34}\}.
\]
Throughout $[-T,T]\times P$, $22<g_4<26$,
$g(\pm T)>10^{-9}$, and
$g_1(-T)<-10^{-6}<0<10^{-6}<g_1(T)$.
There are at most four zeros in the $s$ window, counting multiplicity,
and no root crosses its endpoints. Off the root discriminant there
are exactly $0$, $2$ or $4$ simple real zeros.

If $r_1<\cdots<r_k$ are all distinct stationary points, then $k\leq3$.
When none is a zero of $g$, the root count equals the number of sign
changes in
\[
 g(-T),g(r_1),\ldots,g(r_k),g(T).
\]
Indeed, Rolle's theorem bounds the number of roots and stationary
points, and each intervening interval is strictly monotone. This
criterion remains valid for a degenerate stationary point when its
value is nonzero.

\paragraph{Theorem 2 (complete multiple-root parametrization).}
Set $R_\ell=4T^2=2^{-12}$ and $R_m=R_\nu=32T^3=2^{-16}$.
For every $(s,\ell)\in[-T,T]\times[-R_\ell,R_\ell]$ there is
exactly one $(m,\nu)$ in the closed radius-$(R_m,R_\nu)$ box
satisfying $g=g_1=0$. These solutions form a real analytic graph.
Every multiple root with controls in $P$ belongs to this graph.

For every $|s|\leq T$ there is exactly one control triple
$p_c(s)=(\ell_c,m_c,\nu_c)$ in the radius-$R$ box satisfying
$g=g_1=g_2=0$. It is real analytic and $p_c(0)=0$. Along it,
\[
 16<\frac{d}{ds}g_3(s,p_c(s))<32,\qquad
 0.8s^2<\ell_c(s)<4s^2\quad(s\ne0).
\]
The two branches $s<0$ and $s>0$ consist of ordinary triple zeros
and meet at the unique order-four point in the auxiliary cusp box.
Thus $Q$ is also the unique order-four point in the physical box.

Uniform Banach contractions with fixed dyadic preconditioners prove
both graphs on their entire driver intervals. Their weighted defect
and residual satisfy $q+\eta<1$. With
$J_p=\partial_{(\ell,m,\nu)}(g_0,g_1,g_2)$ and $v=J_p^{-1}e_3$,
\[
 p_c'=-g_3v,\qquad
 \frac{d}{ds}g_3=g_4-g_3\,\partial_p g_3\cdot v.
\]
The verified bounds $-0.24<v_\ell<-0.1$ and the positive last
derivative give the displayed quadratic bounds by integration.
They also show that $\ell_c$ strictly decreases on the negative
branch and increases on the positive branch. All
$0<\ell\leq2^{-17}$ meet both auxiliary branches once, but their
$m,\nu$ values must be checked separately for membership in $P$.
The surface may self-intersect in control space when distinct $s$
values give the same controls; the graph uniqueness does not exclude
that possibility.

\paragraph{Theorem 3 (open regions with all three counts).}
Put $L=2^{-24}$. The open boxes of common control radii
\[
 (L/128,L^2/256,L^2/256)=(2^{-31},2^{-56},2^{-56})
\]
around $(\ell,m,\nu)=(-L,0,0)$, $(-L,-L^2,0)$ and $(L,0,0)$
lie inside $P$ and have respectively $0$, $2$ and $4$ simple real
zeros in $[-T,T]$ for every parameter in the box.

For the first two boxes, the unique minimum of $g_2$ is bracketed
in $|s|\leq2^{-20}$, where $g_2>0$. Thus $g$ is strictly convex
on the full window. Its minimum is positive in the first box;
$g(0)<0$ in the second. For the third box, the uniform signs at
$s=\sqrt L(-3,-1,0,1,3)$ are $(+,-,+,-,+)$. The four-root upper
bound makes the four intervening roots simple and exhaustive.

\paragraph{A finite section.}
At $\ell=L$ both cusp points lie in $P$, with outward bounds
\[
 s_c^-\in[-0.00017264472909,-0.00017264472907],\qquad
 s_c^+\in[0.00017262225524,0.00017262225525].
\]
A simultaneous four-equation contraction establishes a single
control pair with two distinct double zeros, both satisfying $g_2>0$,
at
\[
 s_1\in[-0.00029903062777,-0.00029903025317],\qquad
 s_2\in[0.00029898933163,0.00029898970623].
\]
The four-dimensional Jacobian is nonsingular, so the two critical-value
parameter gradients are independent and the fold branches cross
transversely in this section. The root multiplicities total four;
there are no other roots in the window at those controls. No global
uniqueness claim for all double-fold intersections is made.

\paragraph{A certified region map.}
For display use the explicitly recorded nonsingular dyadic matrix $E$
and the affine chart
\[
 \binom{m}{\nu}=E\binom{L^2(\beta-3/4)}{L^{3/2}\alpha},
 \quad -4\leq\alpha\leq4,\quad-3/2\leq\beta\leq11/2.
\]
Every chart control lies in $P$. This is a known affine map of the
original controls, not an asserted exact quartic normal form.
Of the $64\times56$ closed cells of side $1/8$, $430$, $2184$ and
$190$ have proved counts $0$, $2$ and $4$, respectively.
The other $780$ cells remain unclassified by this grid. Colored
cells have a proof for every parameter in them; no interpolation is
used to assign a count to gray cells.

\paragraph{Verification.}
New majorants through derivative order $60$ include all three
control variations. A total-order-seven expansion in
$sD-\ell D^2/4+mD^4/16-\nu D^6/64$ uses the frozen exact-Q
jets through order $54$; order-eight remainders are explicit.
Separate rational implementations reconstruct the normalized jets,
Taylor bounds, uniform contractions, open witnesses, all $2804$
colored cells and $90$ geometric sample contractions. The latter
comprise two cusps, one double-fold intersection, $81$ fold points
and six simple roots. There are also $27$ direct rigorous integral
comparisons and nine separate two-precision numerical checks.

The rational checkers assume the prior integral enclosures and the
new unnormalized positive majorants. The analytic arguments and
literature priority await external review. All results concern the
fixed modified kernel and the stated finite window; no global zero
count, original fixed-kernel order-four point, or physical stability
interpretation is asserted.
